\section{HBZ/rANS 동역학계와 fixed-input failure exclusion}
\label{sec:rans-frontier}

\subsection{왜 success 사후조건과 종료성은 다른 주장인가}

signature prefix의 역함수(inverse function), HBZ coefficient--symbol leaf,
그리고 actual encoder가 정준 기호열(canonical symbol stream)을 순수 byte
trace로 보내는 간선(edge), exact trace를 actual decoder가 정준 기호열로
되돌리는 간선(edge), actual copy를 사이에 둔 core 합성(composition), 그리고
production SignaturePack/Unpack의 full-HBZ wrapper 경계는 각각 닫혔다.
Week~15는 canonical all-zero HBZ가 만드는 all-6 stream에서 반환하는 actual
encoder의 failure를 배제하고 production zero round trip까지 올렸다
\cite{week15-report,rans-actual-success-witness}. 따라서
\claimid{OBL-RANS-ACTUAL-SUCCESS-WITNESS}는 fixed-input Hoare
부분정확성(partial correctness)으로 \statusproved 다. 다만 반환·실행 존재를
보장하는 termination/losslessness 정리는 아니다. 이 결과 이후 프로젝트는
Week~16 KeyGen/Sign을 거쳐 현재 \CurrentGate 로 이동했다.

rANS 자체의 정수 산술(integer arithmetic)은 복잡하지 않다. 어려움은 encoder가
symbol을 뒤에서부터
읽고 byte buffer cursor를 감소시키는 반면, decoder는 복사된 buffer를 앞에서부터
읽는다는 구현 방향성에 있다. 순수 역함수 정리(pure inverse-function theorem)를
실제 두 중첩 loop의 state와 결합하는 불변식(invariant)이 핵심이다. Week~11은
reverse encoder 쪽 결합을, Week~12는 forward decoder 쪽 결합을, Week~13은
두 정리(theorem) 사이의 actual copy와 전제 운송(premise transport)을 완료했다.
Week~14는 이 core 정리(theorem)를 production
\identifier{_encode_hb_z1_full}/\identifier{_decode_hb_z1_full}의
prepare/apply·frame·size 의미론(semantics)과 합성했다. Week~15는 actual prepare가
all-6 stream을 만든다는 사실, 1020-byte normalization budget, actual failure의
필요조건을 합성해 반환하는 fixed input 실행의 nonzero-size branch를 강제한다.
HBZ에 남은 간극(gap)은 이 실행 자체의 종료·존재, 임의의 canonical 입력에 대한
전역 성공, malformed rejection과 probabilistic losslessness다.

\subsection{13-symbol 분할(partition)}

HBZ 정준 계수(canonical coefficient) \(h_i\in[-6,7)\)를
\(s_i=h_i+6\)으로 옮기면 alphabet은
\[
  \Sigma=\{0,1,\ldots,12\}
\]
이다. rANS scale을 \(M=1024\)라 하고 각 symbol의 누적 시작점과
빈도(frequency)를
\[
\begin{aligned}
(c_s)_{s=0}^{12}
  &=(0,1,2,3,8,66,312,710,957,1016,1021,1022,1023),\\
(f_s)_{s=0}^{12}
  &=(1,1,1,5,58,246,398,247,59,5,1,1,1)
\end{aligned}
\]
로 둔다. certificate는
\[
  f_s>0,\qquad \sum_{s=0}^{12}f_s=M,\qquad
  [c_s,c_s+f_s)\text{가 }[0,M)\text{를 분할함}
\]
을 보인다.

이 분할은 residue slot \(y\bmod M\)에서 symbol을 유일하게 복원하는
유한 분할(finite partition)이다. 실제 encoder/decoder의 packed tables가 정확히
이 \(c_s,f_s\)와
symbol lookup을 나타낸다는 사실도 generated literal arrays에 대해 인증된다.

\subsection{순수 단계는 유클리드 나눗셈(Euclidean division)의
역함수(inverse function)다}

고정 기호(fixed symbol) \(s\)에 대해
\begin{align*}
  E_s(x)
    &=\left\lfloor\frac{x}{f_s}\right\rfloor M
      +c_s+(x\bmod f_s),\\
  D_s(y)
    &=\left\lfloor\frac{y}{M}\right\rfloor f_s
      +(y\bmod M)-c_s
\end{align*}
로 둔다. 유클리드 나눗셈(Euclidean division)
\(x=f_s\lfloor x/f_s\rfloor+(x\bmod f_s)\)와
\(0\le x\bmod f_s<f_s\)에서 즉시
\[
  D_s(E_s(x))=x
\]
를 기대할 수 있다. EasyCrypt의
\theoremname{pure_rans_step_quotient},
\theoremname{pure_rans_step_slot},
\theoremname{pure_rans_step_inverse}가 이 계산을 분리해 증명한다.

실제 encoder는 나눗셈 대신 table의 reciprocal, shift, bias를 사용하는 fast
formula를 쓴다. 별도 certificate는
\[
  1\le x<x_{\max}(s),\qquad
  x_{\max}(s)=2^{21}f_s
\]
에서 fast quotient가 유클리드 몫(Euclidean quotient)과 같고, 따라서 fast step도
\(D_s\)에 의해 역산됨을 보인다. 이것은 단순히 ``표 값이 맞아 보인다''는
계산 검사가 아니라 실제 word field의 정수 의미를 연결한 정리(theorem)다.

\begin{verified}[순수 및 fast step]
정준 기호(canonical symbol) \(0\le s<13\)과 명시된 상태 범위(state range)에서
\[
  D_s(E_s(x))=x,\qquad
  D_s(E^{\mathrm{fast}}_s(x))=x
\]
가 컴파일된다. 전제 \(s=0,x=1\)의 concrete witness도 있어 정리(theorem)는 공허하지
않다.
\end{verified}

\subsection{정규화 바이트(normalization byte)는 잃어버린 몫(quotient)의
좌표(coordinate)다}

state가 \(x_{\max}(s)\)보다 크면 fast step을 적용하기 전에 byte radix
\(B=256\)로 0, 1, 2회 나누어 state를 줄인다. 순수 정의는
\[
  \ell_s(x)\in\{0,1,2\},\qquad
  r_s(x)=\left\lfloor\frac{x}{B^{\ell_s(x)}}\right\rfloor
\]
이고, 나눌 때 버린 낮은 byte들을 decoder가 읽을 순서의 list
\(\beta_s(x)\)로 보존한다.

핵심 readback 명제(proposition)는
\[
  \operatorname{ReadBytes}(r_s(x),\beta_s(x))=x.
\]
한 byte일 때는 \(x=B\lfloor x/B\rfloor+(x\bmod B)\)이고, 두 byte일 때는
이 항등식(identity)을 두 번 적용한다. 실제 encoder는 낮은 byte를 먼저 쓰면서 cursor를
감소시키므로 최종 forward segment에서 두 byte의 순서는
\([\mathsf{high},\mathsf{low}]\)다. 이 방향 전환을 놓치면 순수 정리(theorem)는 맞아도
actual decoder와 맞지 않는다.

\sourcepath{Mode2RansNormalization.ec}는 실제 \(W32\) shift/truncate/or와
\(W64\) cursor decrement가 이 정수 연산과 같은 의미를 갖는다는
leaf 정리(theorem)를
제공한다.

\subsection{전체 symbol list의 순수 trace}

symbol list를 \(s_0,\ldots,s_{n-1}\), \(n=1024\)라 하고 마지막 state를
\(x_n=L=2^{23}\)로 둔다. encoder가 뒤에서 앞으로 진행하도록
\[
\begin{aligned}
  \ell_j&=\ell_{s_j}(x_{j+1}),\\
  r_j&=\left\lfloor x_{j+1}/256^{\ell_j}\right\rfloor,\\
  x_j&=E^{\mathrm{fast}}_{s_j}(r_j),\\
  \beta_j&=\beta_{s_j}(x_{j+1})
\end{aligned}
\]
로 재귀적으로(recursively) 정의한다. decoder-facing byte string은
\[
  \mathcal B=\operatorname{LE32}(x_0)
    \Vert\beta_0\Vert\cdots\Vert\beta_{n-1}
\]
이다. cursor cut은
\[
  c_0=4,\qquad c_{j+1}=c_j+|\beta_j|
\]
로 둔다.

프로젝트의 이름 \(xs/cuts/bytes\)는 각각
\((x_0,\ldots,x_n)\), \((c_0,\ldots,c_n)\), \(\mathcal B\)를 가리킨다. 이 세
대상은 \sourcepath{Mode2RansByteStack.ec}에 정의되어 있다. 그 파일의
\identifier{valid_rans_trace}는 임의의 witness를 전제로 받지 않고 canonical
symbol list로부터 이들을 직접 구성한다.

\begin{verified}[트레이스(trace)의 귀납적 역산(inductive inversion)]
canonical list에 대해
\[
  2^{23}\le x_j<2^{31}
\]
가 모든 \(j\)에서 성립한다. 또한 \(x_j\)의 residue slot은 \(s_j\)를 선택하고,
\(D_{s_j}(x_j)=r_j\)이며, \(\beta_j\)를 읽으면 \(x_{j+1}\)가 복원된다.
\end{verified}

따라서 순수 decoder의 리스트 수준 귀납증명(list-level inductive proof)에 필요한
수학적 내용은 준비되어 있다. actual encoder state가 이 trace 좌표를
나타낸다는 refinement는 Week~11에, actual decoder state가 같은 좌표를
순방향(forward)으로 소비한다는 refinement는 Week~12에 닫혔다.

\subsection{Week 10에서 감사된 여섯 encoder 타깃}

Week~10은 필요한 바깥 루프 불변식(outer-loop invariant)을 서로 다른 운송
문제로 분해한 다음 \WeekTenTargetCount 개 파일을 authored target manifest에
편입했다. 여섯 파일은 모두 \identifier{-no-eco}로 fresh compile되지만, 각
국소 명제(local proposition)의 완료와 부모 refinement의 완료는 구분해야 한다.

\begin{itemize}
  \item \sourcepath{Mode2RansArrayListBridge.ec}: actual byte array의 symbol
        stream을 순수 list와 suffix로 옮기고, segment 일치와 prefix frame을
        표현한다. 배열(array)--리스트(list) 다리와 한 기호 trace
        점화식(recurrence)은 \statusproved 다.
  \item \sourcepath{Mode2RansEncoderWordStep.ec}: packed reciprocal, shift,
        complement를 쓰는 실제 word 식을 순수 fast step의 정수 의미와
        연결한다. literal mode~2 table에 대한 word-level 등식(equality)은
        \statusproved 다.
  \item \sourcepath{Mode2RansEncoderInnerProgress.ec}: normalization 내부
        loop의 순수한 0/1/2회 진행을 state, cursor, 이미 쓴 byte segment로
        좌표화한다. 이 순수 보조정리(pure lemma)는 \statusproved 다.
  \item \sourcepath{Mode2RansEncoderSerialization.ec}: 마지막 \(W32\) state의
        little-endian 네 바이트, 바깥 구간 frame, 성공 시 size 산술을 분리한다.
        이 leaf 명제(leaf proposition)는 \statusproved 다.
  \item \sourcepath{Mode2RansEncoderTrace.ec}: actual
        내부 loop가 사용하는 유한 phase, byte prepend, segment, prefix frame,
        no-wrap 보조정리(lemma)를 제공한다.
  \item \sourcepath{Mode2RansEncoderActualInner.ec}: 실제 generated
        \identifier{_rans_encode} 내부 loop에 위 불변식(invariant)을 설치하고,
        별도 실제 Hoare 정리(theorem)로 성공 offset과 size 범위(range)를
        산출한다.
\end{itemize}

\begin{caution}[Week 10의 국소 \statusproved 와 당시 부모 \statuspartial]
Week~10 기준선에서 위 파일들은 보조정리(lemma)와 국소 실제
부분정확성(actual partial correctness)만 인증했다. 당시에는 전체 actual suffix와
순수 \identifier{trace_bytes}의 등식(equality)이 없어
\claimid{OBL-RANS-ENCODE-REFINEMENT}가 \statuspartial 이었다. 바로 그 부모
간선(edge)을 Week~11의 다음 절이 닫는다. 성공 분기(success branch)의
도달가능성(reachability)과 종료성(termination)은 이 Week~10 타깃들이 인증하지
않았다. fixed all-6 failure exclusion은 Week~15의 별도 정리이고, 종료성은 현재도
열려 있다.
\end{caution}

\subsection{Week 11에서 닫힌 actual encoder 불변식(invariant)}

실제 reverse encoder의 바깥 루프 index를 \(i\), 현재 word state를 \(x\), 뒤에서
감소하는 buffer cursor를 \(off\), buffer를 \(enc\)라 하자. 순수 trace
\[
  T_i=\operatorname{EncodeTrace}
       (\operatorname{SymbolSuffix}(symbols_0,i))=(x_i,b_i)
\]
를 두면 성공 가지의 live 불변식(live invariant)은 다음 다섯 조건이다
\cite{rans-encoder-invariant}.
\[
\begin{gathered}
  0\le i\le1024,\qquad x=x_i,\qquad off+|b_i|=1024,\\
  \operatorname{SegmentMatches}(enc,off,b_i),\qquad
  \operatorname{PrefixFrame}(enc_0,enc,off),\qquad 4\le off.
\end{gathered}
\]

Week~11에는 \WeekElevenTargetCount 개 타깃이 더해졌다. tail-aware inner
불변식(invariant)은 최초 배열 \(enc_0\)와
\(\operatorname{EncodeTrace}(tail).bytes\)를 actual byte write 전후에 보존하고,
실제 guard 종료에서 정확한 \(0..2\) 정규화 길이(normalization length)를
회복한다. generated word-step 보조정리(lemma)는 실제 table load와 protect
erase 뒤의 word update를 순수 fast step과 같게 만들고, 마지막 네 generated
store는 직렬화된 상태(serialized state)를 누적 suffix 앞에 붙인다.

그 결과 실제 절차의 성공 사후조건(success postcondition)은
\[
\begin{gathered}
  \operatorname{SegmentMatches}
  \bigl(enc',off',\operatorname{TraceBytes}
  (\operatorname{SymbolList}(symbols_0))\bigr),\\
  off'+\left|\operatorname{TraceBytes}
  (\operatorname{SymbolList}(symbols_0))\right|=1024,\\
  \operatorname{PrefixFrame}(enc_0,enc',off')
\end{gathered}
\]
로 컴파일되었다.

\begin{verified}[\claimid{OBL-RANS-ENCODE-REFINEMENT}]
\theoremname{actual_rans_encode_trace_closure}는 실제 generated
\identifier{_rans_encode}가 반환한다면 실패하거나, 성공 시 위 세 관계(relation),
offset/size 범위(range), \(bad=0\)을 모두 보인다.
\theoremname{actual_rans_encode_trace_refinement}는 전체 반환
접미부(returned suffix) 등식(equality)을 부모 refinement 표면에 명시적으로
노출한다. 성공은 전제가 아니라 actual 반환값에 따른 결론의 한 가지(branch)다.
\end{verified}

\begin{caution}[부분정확성(partial correctness)의 한계]
이 정리(theorem)는 actual encoder의 종료성(termination), 성공
도달가능성(success reachability), all-6 symbol 성공 증인(success witness)을
그 자체로 보이지 않는다. Week~11 당시
\claimid{OBL-RANS-ACTUAL-SUCCESS-WITNESS}는 \statuspartial 이었고, Week~15의
별도 정리가 fixed all-6 반환 실행의 failure만 배제한다.
\end{caution}

\subsection{Week 12에서 닫힌 actual decoder 불변식(invariant)}

Week~12는 actual forward decoder index를 \(i\), state를 \(x\), 읽기 cursor를
\(off_d\)라 할 때 다음 불변식(invariant)을 generated outer/inner loop에 직접
설치했다 \cite{week12-report,rans-decoder-invariant}.

\begin{enumerate}
  \item 이미 출력한 symbol prefix가 \(s_0,\ldots,s_{i-1}\)와 같다;
  \item 현재 state가 \(x_i\)이고 현재 cursor가 \(c_i\)다;
  \item table lookup이 \(s_i\)를 선택하고, normalization segment를 정확히
        \([c_i,c_{i+1})\)에서 읽는다;
  \item 출력 배열의 아직 쓰지 않은 tail에 대한 frame이 유지된다;
  \item 종료 시 \(x=2^{23}\), \(off_d=size\), \(bad=0\)이다.
\end{enumerate}

역사적 \theoremname{actual_rans_decode_mode2_control}은 실제 tables/state fields와
\(bad\in\{0,1\}\)만 고정했다. 새 입력 술어(input predicate)
\identifier{actual_mode2_decoder_trace_input}은 expected symbol array, exact
\identifier{trace_bytes}, 모든 normalization read의 pointwise relation, state
fields \((1024,encoded\_size,13)\), 실제 generated tables를 한데 묶는다.

\begin{verified}[\claimid{OBL-RANS-DECODE-REFINEMENT}]
\theoremname{actual_rans_decode_trace_refinement}는 실제 generated
\identifier{_rans_decode}를 직접 열어 위 exact-trace 입력 아래 모든 1024개
symbol을 복원하고, \(off=encoded\_size\), \(bad=0\), decoded tail frame을
결론낸다. \theoremname{decoder_outer_trace_exit_components}는 내부적으로
\(x=2^{23}\)을 유도하여 actual final-state reject도 통과시킨다.
\end{verified}

이 정리(theorem)는 Hoare 부분정확성(Hoare partial correctness)이므로 decoder
종료성(termination)을 보이지 않는다. exact-trace 입력은 비순환적
(non-circular)이며, Week~13은 actual harness가 그 입력을 산출한다는 별도 합성
정리(composition theorem)를 추가했다.

\subsection{Week 13에서 닫힌 actual harness core 합성(composition)}

\sourcepath{Mode2RansActualInverse.ec}의 harness는 추상 조립도가 아니라 세 실제
generated 절차를 다음 순서로 호출한다. 같은 파일의 정리(theorem)는 이 호출의
control만 닫고, Week~13의 의미론적 closure는
\sourcepath{Mode2RansCoreActualInverse.ec}에 따로 있다.
\[
\mathsf{RansEncode}
\longrightarrow
\mathsf{CopyEncodedSuffix}
\longrightarrow
\mathsf{RansDecode}.
\]
단, encoder가 \(bad=0\)을 반환한 경우에만 뒤의 두 호출을 실행한다.

기존 harness control 결론은
\[
\begin{aligned}
  &\mathit{encoderBad}\in\{0,1\},\\
  &\mathit{decoderRan}\Longleftrightarrow\mathit{encoderBad}=0,\\
  &\mathit{decoderRan}\Longrightarrow
    (count,size,m)=(1024,\mathit{encodedSize},13)
\end{aligned}
\]
이다. Week~13의 \sourcepath{Mode2RansCoreCompositionBridge.ec}는 encoder
\identifier{segment_matches}와 copy \identifier{slice_eq}를 copied buffer의
exact trace로 합성하고, 이를 decoder의 pointwise read와 configured state 입력
술어(input predicate)로 운송한다. 특히 modular (W64) 뺄셈을 정수 뺄셈(integer
subtraction)으로 몰래 바꾸지 않고 actual success bound에서 no-wrap을 유도한다.

\begin{verified}[\claimid{OBL-RANS-CORE-INVERSE}]
\theoremname{actual_rans_encode_copy_decode_inverse}는 actual encoder 성공을
전제로 받지 않는다. 실패 시 actual copy와 decoder가 실행되지 않고 초기 decoded
array/state가 보존된다. 성공 시 다음 의미론적 결론(semantic conclusion)이 같은
harness 정리(theorem)의 사후조건(postcondition)에 들어간다.
\[
  \mathit{decoderBad}=0,\quad
  off_{\mathrm{final}}=size,\quad
  decoded[0..1024)=symbols[0..1024).
\]
decoder 하위 증명은 내부 종료 상태(internal exit state)
\(x=2^{23}\)도 유도하지만 harness ABI는 이를 공개하지 않는다. decoded tail
\([1024,2048)\)은 frame으로 보존된다. 이 결과는 성공 조건부
부분정확성(success-conditioned partial correctness)이다.
\end{verified}

Week~13 종료 시점에는 실제 success branch를 강제하는 EasyCrypt 정리가 없었다.
pure trace의 concrete witness, canonical all-zero precondition witness 또는
decoder state의 size-4 witness는 concrete all-6 입력에서 actual encoder가
\(bad=0\)을 반환한다는 사후조건을 대신하지 않았다. Week~15는 바로 이 간극을
budget과 actual failure-cause의 모순으로 닫는다.

\subsection{Week 14에서 닫힌 full HBZ wrapper의 올바른 판독}

actual encoder는 실패를 표현할 수 있으므로, canonical \(h\)만으로 무조건적인
round trip을 선언해서는 안 된다. Week~14에서 증명된 부모 명제(parent proposition)는
\[
  size=0\quad\lor\quad
  \bigl(size\ne0\land bad=0\land
        decoded[0..1024)=h[0..1024)\bigr)
\]
처럼 failure branch를 드러내야 한다. nonzero size bounds와 buffer frames도
함께 필요하다.

\begin{verified}[\claimid{OBL-SIG-HBZ-ENCODE-DECODE}]
\theoremname{actual_hbz_full_encode_decode_inverse_mode2}가 focused actual
full-HBZ harness에서 위 failure/success 합(disjunction)을 증명한다.
\theoremname{signature_pack_unpack_hbz_full_actual_exact}는 production
SignaturePack/Unpack harness가 focused harness와 반환 tuple 및
\identifier{decoder_ran} flag에서 exact함을 보이고,
\theoremname{signature_pack_unpack_hbz_full_inverse_mode2}가 같은 계약을
production 경계에 올린다. 성공 가지에는 \(4\le size\le1024\), coefficient
prefix 복원, coefficient/suffix frame과 prepared-symbol trace witness가 포함된다.
\end{verified}

\subsection{Week 15의 fixed all-6 budget과 actual failure 배제}

\identifier{zero_hbz}의 실제 32-bit load가 모두 0임을 byte projection까지 열어
증명하면 actual prepare가 \(bad=0\)과 1024개 symbol~6 prefix를 반환한다. literal
table에서 \(f_6=398\), \(x_{\max}(6)=2^{21}\cdot398\)이고,
\theoremname{symbol6_normalization_len_le1}는 \(2^{23}\le x<2^{31}\)인 각
symbol~6 step이 normalization byte를 최대 하나만 방출함을 준다.

초기 state에서 계산한 첫 네 상태는
\[
  8388608\to21582496\to55528910\to142868116\to367580518
\]
이고 네 step 모두 0-byte다. 따라서 나머지 1020개 step에 각 한 byte를 허용해도
normalization length는 최대 1020이며, 마지막 state 네 바이트를 더한 total
trace는 \([4,1024]\)에 들어간다. 이는 관찰된 정확한 174/178-byte 실행을
정리로 가장하지 않는 보수적 capacity argument다.

\begin{verified}[\claimid{OBL-RANS-ACTUAL-SUCCESS-WITNESS}]
\theoremname{actual_rans_encode_failure_trace_cause}는 actual encoder의 nonzero
\(bad\)가 trace length \(>1020\)을 함의한다고 말한다. 위 all-6 상계와
모순시키면 \theoremname{actual_rans_encode_all_six_success}가 반환하는 actual
core 실행마다 \(bad=0\), size/frame/exact-trace 사후조건을 준다.
\theoremname{actual_encode_hb_z1_full_zero_success},
\theoremname{signature_pack_hbz_zero_success_mode2},
\theoremname{signature_pack_unpack_hbz_zero_success_mode2}가 이 결론을 실제
prepare를 포함한 focused full wrapper와 production SignaturePack/Unpack zero
round trip으로 올린다.
\end{verified}

\begin{caution}[\identifier{GO-WITNESS}는 \identifier{GO-WITNESS-LOSSLESS}가 아니다]
위 정리는 모든 \emph{반환하는} fixed-input 실행에서 failure를 배제하는 Hoare
부분정확성이다. fixed-input termination, 실행 존재(non-vacuity), losslessness,
probability-one success 또는 \identifier{phoare}는 증명하지 않았다. \(h\)
codec은 성공 증인이 미완이라서가 아니라 Week~16 MINCORE KeyGen 범위로 전환되어
\statusdeferred 되었다.
\end{caution}

\begin{center}
\begin{tabularx}{\textwidth}{Y L{31mm}}
\toprule
rANS 층(layer) & 상태 \\
\midrule
13-symbol interval/frequency와 actual table certificate & \statusproved \\
pure quotient/slot/fast-step inverse & \statusproved \\
0/1/2-byte normalization과 W32/W64 leaf semantics & \statusproved \\
canonical \(xs/cuts/bytes\) trace와 state bounds & \statusproved \\
actual array/list suffix와 word-step bridge & \statusproved \\
actual 네-byte 최종 serialization 합성 & \statusproved \\
actual encoder 성공 offset/size 범위 & \statusproved\ (부분정확성(partial correctness)) \\
actual 내부 normalization byte/frame/no-underflow &
\statusproved\ (tail-aware exact \(0..2\) exit) \\
actual encoded suffix copy와 tail frame & \statusproved \\
actual encoder--copy--decoder branch control & \statusproved\ (control만) \\
actual encoder가 pure trace 전체를 산출함 &
\statusproved\ (성공 조건부 부분정확성(success-conditioned partial correctness)) \\
actual decoder가 exact trace를 소비함 &
\statusproved\ (1024-symbol 복원, exact consumption, tail frame) \\
actual encoder--copy--decoder core inverse &
\statusproved\ (성공 조건부 부분정확성(success-conditioned partial correctness)) \\
full HBZ encode/decode &
\statusproved\ (production wrapper의 성공 조건부 부분정확성) \\
actual encoder success witness &
\statusproved\ (fixed all-6 Hoare failure exclusion; termination/losslessness 없음) \\
\bottomrule
\end{tabularx}
\end{center}
